\documentclass[12pt]{article}

\usepackage{sbc-template}
\usepackage{graphicx,url}
\usepackage[utf8]{inputenc}
\usepackage[english]{babel}

\sloppy

\title{Relatório Prático: Kernel Mínimo em RISC-V\\
\large Boot, Comunicação Serial e Kernel Panic}

\author{João Pedro Sant Ana de Carvalho}

\address{Faculdade de Computação -- Universidade Federal de Mato Grosso do Sul (UFMS)\\
  Campo Grande -- MS -- Brasil
  \email{santc.joao@gmail.com}
}

\begin{document}

\maketitle

\begin{abstract}
This report documents the implementation of a minimal 32-bit RISC-V operating system kernel for the Operating Systems course. The completed tasks include kernel boot with a dedicated linker script (Task 1), serial communication via OpenSBI calls (Task 2), and a kernel panic mechanism for error handling (Task 4). Tasks 3 (full freestanding C library) and 5 (trap handling) were not implemented in this version. The implemented code was verified on the QEMU emulator.
\end{abstract}

\begin{resumo}
Este relatório apresenta o desenvolvimento de um kernel mínimo para a arquitetura RISC-V de 32 bits, realizado como atividade prática na disciplina de Sistemas Operacionais. O texto foca nas etapas que foram concluídas: o boot do kernel com linker script próprio (Trabalho 01), a saída serial via chamadas ao OpenSBI (Trabalho 02) e a implementação de kernel panic para lidar com falhas (Trabalho 04). Fica registrado que os Trabalhos 03 (biblioteca C básica completa) e 05 (tratamento de exceções) não foram realizados nesta entrega. O funcionamento dos módulos concluídos foi testado no emulador QEMU.
\end{resumo}

\section{Introdução}

O objetivo deste projeto prático é criar um kernel básico do zero (bare-metal) para a arquitetura RISC-V de 32 bits (RV32), rodando no emulador QEMU junto com o firmware OpenSBI.

Das cinco tarefas propostas na disciplina, este relatório aborda as que foram finalizadas e testadas no repositório:
\begin{itemize}
    \item \textbf{Trabalho 01}: Boot do kernel e linker script;
    \item \textbf{Trabalho 02}: Hello World com chamadas ao OpenSBI;
    \item \textbf{Trabalho 04}: Mecanismo de Kernel Panic.
\end{itemize}

Por restrições de tempo e escopo, os \textbf{Trabalhos 03} (biblioteca C com funções de string e memória) e \textbf{05} (tratador de exceções da CPU) \textbf{não foram realizados} nesta entrega.

\section{Atividades Realizadas}

\subsection{Trabalho 01: Boot do Kernel e Linker Script}

A primeira etapa foi definir como o kernel é organizado na memória por meio do linker script (\texttt{kernel.ld}). O endereço base foi colocado em \texttt{0x80200000}, que é o ponto padrão para onde o OpenSBI pula depois de inicializar o hardware.

O linker script divide a imagem nas seções usuais: \texttt{.text}, \texttt{.rodata}, \texttt{.data} e \texttt{.bss}. A seção \texttt{.text.boot} foi colocada logo no começo usando \texttt{KEEP} para garantir que o código de boot fique no endereço inicial. Logo após a seção \texttt{.bss} (delimitada pelos símbolos \texttt{\_\_bss} e \texttt{\_\_bss\_end}), foram reservados 128 KB para a pilha do kernel, com o símbolo \texttt{\_\_stack\_top} marcando o topo da pilha.

No arquivo \texttt{kernel.c}, a função \texttt{boot} foi escrita em assembly inline com o atributo \texttt{naked}, evitando que o compilador adicione instruções extras de início ou fim de função. Essa função apenas ajusta o registrador de pilha (\texttt{sp}) para \texttt{\_\_stack\_top} e pula para \texttt{kernel\_main}:

\begin{verbatim}
__attribute__((section(".text.boot")))
__attribute__((naked)) void boot(void) {
    __asm__ __volatile__(
        "mv sp, %[stack_top]\n"
        "j kernel_main\n"
        : : [stack_top] "r"(__stack_top)
    );
}
\end{verbatim}

Ao entrar em \texttt{kernel\_main}, a primeira ação é limpar a seção \texttt{.bss}, preenchendo todos os seus bytes com zero através de um \texttt{memset} local. Isso garante que variáveis globais e estáticas não comecem com valores aleatórios da memória.

\subsection{Trabalho 02: Saída Serial com OpenSBI}

Para conseguir imprimir mensagens no terminal do QEMU sem drivers complexos, utilizou-se a Supervisor Binary Interface (SBI) fornecida pelo OpenSBI através da instrução \texttt{ecall}.

A função \texttt{sbi\_call} organiza os parâmetros nos registradores da CPU: \texttt{a7} recebe o Extension ID (EID), \texttt{a6} o Function ID (FID), \texttt{a0} a \texttt{a5} levam os argumentos, e o retorno vem em \texttt{a0} e \texttt{a1} (guardados na estrutura \texttt{struct sbiret}). Com isso, foi implementada a função \texttt{putchar} usando a extensão clássica \textit{Console Putchar} (EID 1, FID 0):

\begin{verbatim}
void putchar(char ch) {
    sbi_call(ch, 0, 0, 0, 0, 0, 0, 1);
}
\end{verbatim}

A partir do \texttt{putchar}, foi criada uma rotina simples para percorrer strings terminadas em \texttt{'\textbackslash0'} e enviar caractere por caractere para a serial, permitindo imprimir mensagens como \texttt{"Hello World!"}.

\subsection{Trabalho 04: Mecanismo de Kernel Panic}

O quarto trabalho teve como foco criar uma forma simples de interromper o sistema quando acontece um erro fatal irrecuperável. No arquivo \texttt{kernel.h}, foi definida a macro \texttt{PANIC}:

\begin{verbatim}
#define PANIC(fmt, ...) \
    do { \
        printf("PANIC: %s:%d: " fmt "\n", \
               __FILE__, __LINE__, ##__VA_ARGS__); \
        while (1) {} \
    } while (0)
\end{verbatim}

A macro usa o padrão \texttt{do \{ ... \} while(0)} para funcionar bem dentro de condicionais (\texttt{if/else}). Foram utilizadas as macros do compilador \texttt{\_\_FILE\_\_} e \texttt{\_\_LINE\_\_} para mostrar no terminal o arquivo e a linha onde a falha aconteceu, além da mensagem passada como parâmetro. Depois de imprimir, ela entra em um loop infinito (\texttt{while(1)}), travando a CPU e impedindo que o kernel continue executando instruções com o estado corrompido.

\section{Atividades Não Realizadas}

Para fins de transparência sobre o escopo entregue, detalha-se o que ficou pendente:

\begin{itemize}
    \item \textbf{Trabalho 03 (Biblioteca C Básica):} A proposta original pedia a implementação de funções de manipulação de memória e strings (\texttt{memcpy}, \texttt{strcpy}, \texttt{strcmp}), além de tipos próprios. No código atual, foi implementado apenas um \texttt{memset} enxuto direto no \texttt{kernel.c} e um \texttt{printf} básico em \texttt{common.c} para viabilizar a impressão do pânico. A biblioteca completa planejada em \texttt{common.c} e \texttt{common.h} não foi finalizada.
    \item \textbf{Trabalho 05 (Tratamento de Exceções):} Esta etapa previa a configuração do registrador \texttt{stvec}, a rotina \texttt{kernel\_entry} em assembly para salvar todos os registradores na pilha (estrutura \texttt{trap\_frame}), e a leitura dos registradores CSR (\texttt{scause}, \texttt{stval}, \texttt{sepc}) para identificar erros de hardware e instruções ilegais. Esse tratador não foi implementado, portanto o kernel ainda não captura exceções da CPU.
\end{itemize}

\section{Resultados e Validação}

O kernel foi compilado com o Clang/LLD para RISC-V 32-bit em modo \textit{freestanding} e testado no QEMU através do script \texttt{run.sh}.

Os testes confirmaram o funcionamento das partes implementadas:
\begin{enumerate}
    \item \textbf{Boot e Pilha:} A inspeção dos símbolos gerados no \texttt{kernel.map} e no desassemblador confirmou que o ponto de entrada \texttt{boot} foi colocado exatamente em \texttt{0x80200000} e a pilha começou em \texttt{0x80220018}.
    \item \textbf{Saída Serial e Kernel Panic:} No \texttt{kernel\_main}, executou-se intencionalmente a chamada \texttt{PANIC("booted!")} para testar a parada de emergência:
\begin{verbatim}
PANIC: kernel.c:42: booted!
\end{verbatim}
    A mensagem apareceu no terminal indicando corretamente o arquivo e a linha, e a linha seguinte (\texttt{printf("unreachable here!\textbackslash n")}) não chegou a ser executada, confirmando que a CPU realmente travou no loop infinito.
\end{enumerate}

\section{Conclusão}

Neste trabalho, foi possível colocar um kernel RISC-V para rodar do zero em bare-metal, estabelecer a comunicação serial com o OpenSBI para exibição no terminal e criar um mecanismo de kernel panic para contenção de erros. Os Trabalhos 03 e 05 não foram concluídos nesta entrega e ficam como etapas necessárias antes de avançar para tópicos como paginação de memória e escalonamento de processos.

\begin{thebibliography}{9}
\bibitem{nuta2023}
Nuta, S. (2023). \textit{Operating System in 1,000 Lines}.

\bibitem{sbispec}
RISC-V SBI Task Group (2022). \textit{RISC-V Supervisor Binary Interface Specification}.
\end{thebibliography}

\end{document}
